% !TeX program = tectonic
\documentclass[11pt,a4paper]{article}
\usepackage{fontspec}
\usepackage[english]{babel}
\babelfont{rm}[Language=Default]{Liberation Serif}
\babelfont[Language=Default]{sf}{Carlito}
\babelfont[Language=Default]{tt}{DejaVu Sans Mono}
\usepackage{amsmath,amssymb,amsthm}
\usepackage{geometry}
\geometry{a4paper, top=2.4cm, bottom=2.4cm, left=2.4cm, right=2.4cm}
\usepackage{booktabs,tabularx,enumitem,xcolor,tcolorbox}
\tcbuselibrary{breakable,skins}
\definecolor{linkblue}{HTML}{1F4E79}
\definecolor{statgreen}{HTML}{2F6B3A}
\definecolor{goldacc}{HTML}{8B7E5A}
\usepackage[colorlinks=true,linkcolor=linkblue,urlcolor=linkblue,unicode=true]{hyperref}
\theoremstyle{plain}
\newtheorem{theorem}{Theorem}
\newtheorem{lemma}[theorem]{Lemma}
\newtheorem{proposition}[theorem]{Proposition}
\newtheorem{corollary}[theorem]{Corollary}
\theoremstyle{definition}
\newtheorem{definition}[theorem]{Definition}
\theoremstyle{remark}
\newtheorem{remark}[theorem]{Remark}
\newtcolorbox{statusbox}[1][]{colback=statgreen!5,colframe=statgreen!75,
  title={\textbf{Summary of results:} #1},fonttitle=\small,boxrule=0.7pt,arc=2pt,breakable}
\newtcolorbox{databox}[1][]{colback=goldacc!6,colframe=goldacc!80,
  title={\textbf{Protocol data:} #1},fonttitle=\small,boxrule=0.7pt,arc=2pt,breakable}
\newtcolorbox{verifybox}[1][]{colback=linkblue!5,colframe=linkblue!80,
  title={\textbf{Verification:} #1},fonttitle=\small,boxrule=0.7pt,arc=2pt,breakable}
\widowpenalty=10000
\clubpenalty=10000
\setlength{\parskip}{0.35em}
\newcommand{\CC}{\mathbb{C}}
\newcommand{\RR}{\mathbb{R}}
\newcommand{\ZZ}{\mathbb{Z}}
\newcommand{\QQ}{\mathbb{Q}}
\newcommand{\Ga}{\Gamma}
\newcommand{\Bfun}{\mathrm{B}}
\newcommand{\im}{\operatorname{Im}}
\newcommand{\re}{\operatorname{Re}}
\newcommand{\lcm}{\operatorname{lcm}}
\newcommand{\SNF}{\operatorname{SNF}}
\newcommand{\Jac}{\operatorname{Jac}}
\newcommand{\rank}{\operatorname{rank}}
\newcommand{\Ind}{\operatorname{Index}}
\newcommand{\disc}{\operatorname{disc}}
\begin{document}
\begin{center}
{\LARGE\bfseries Certificate E: Flow Termination}\\[0.4em]
{\large the exact time t* = lcm and the terminal cycle}\\[0.6em]
{\normalsize Isaev Iskhak Khamzatovich}\\[0.2em]
{\normalsize The <<Dynamic Principle>> Program --- the \texttt{hodge-laboratory}}\\[0.15em]
{\small Standalone theorem monograph · Монография 13/16}
\end{center}
\vspace{0.4em}
{\color{linkblue}\hrule height 0.8pt}
\vspace{0.8em}

\section{Setting}

Certificate E proves the termination of the discrete pipeline flow
and determines the exact termination time. Every laboratory
experiment provably stops --- a fundamental stability property.

\begin{theorem}[E1--E4: termination]\label{th:Estand}
\begin{enumerate}[leftmargin=2em]
\item[(E1)] The state space of the flow is finite.
\item[(E2)] A discovery monovariant exists: every iteration either
opens a new state or moves the flow to the terminal regime.
\item[(E3)] By the pigeonhole principle the flow terminates; the
terminal state is a cycle.
\item[(E4)] The exact termination time is
$t^*=\lcm\bigl(W/\gcd(a,W),\,H/\gcd(b,H)\bigr)$.
\end{enumerate}
\end{theorem}

\begin{proof}
(E1) States are (position, phase) pairs on a finite lattice with a
finite phase group: at most $4WH$ states.
(E2) The monovariant is the set of visited cells: strictly growing
in the discovery regime, bounded by the grid size.
(E3) Strict growth of a finite monovariant stops; the stop is a
cyclic repetition of the configuration by the definition of the
transition.
(E4) The horizontal return is after $W/\gcd(a,W)$ steps, the
vertical after $H/\gcd(b,H)$; the joint return is the LCM.
Substituting the reference ($48\times48$, step $(1,1)$) gives
$t^*=48$ --- and the flow terminates in exactly $48$ steps with the
full coordinate cross-check $48/212/432/114$.
\end{proof}

\begin{databox}[the eight protocol points]
E-1: scanning $64$ step pairs --- the monovariant grows, the flow
terminates. E-2: the $48\times48$ pipeline: $t^*=48$; the
$48/212/432/114$ cross-check complete. E-3: the $(7,22)$ correction
step --- termination under the correction rules. E-4: the terminal
cycle of length $4$ --- the $\mu_4$-orbit of the step. E-5: scales
$48\to384$ --- the $t^*$ formula persists. E-6: the K3 layer
(reflections with braking) terminates. E-7: the Klein layer:
Singer generation of order $7$. E-8: the Klein layer:
$\mu_4$-orbits of length $4$. Verdict: $8/8$ PASS.
\end{databox}

\begin{remark}[a cycle --- a theorem, not a hypothesis]
The cyclic structure of the terminal state follows from finiteness
and step periodicity; the cycle length $4$ is exactly the
$\mu_4$-orbit. This is a theorem about the \emph{structure} of the
terminal regime, not merely its existence.
\end{remark}

\begin{statusbox}[summary]
Certificate E accepted ($8/8$): the flow terminates, the
termination time is computed exactly, the terminal structure is
described.
\end{statusbox}

\begin{verifybox}[reproduction]
\texttt{python3 laboratory.py} $\to$ ``Certificates $\to$ E''; the
designer: an arbitrary $W\times H$ grid and step $(a,b)$ --- the
lab prints $t^*$ and verifies by enumeration.
\end{verifybox}

\vfill
{\color{linkblue}\hrule height 0.6pt}
\vspace{0.5em}
{\small\color{gray}
License: individual exclusive license of Isaev Iskhak Khamzatovich.
All rights to the computations, formulas, and texts belong to the author.
Verification: \texttt{python3 laboratory.py} (``Stands''/``Certificates''),
Lean: \texttt{verification/lean/HodgeLaboratory.lean}.
}
\end{document}
